\section*{Chapter \ref{ch:g11:stat} --- Descriptive Statistics}

\begin{solution}{exo:g11:stat:1}
$n = 8$, sum $71$: \omterm{def:g11:stat:mean}{mean} $\bar x = \frac{71}{8} = 8.875$. \omterm{def:g11:stat:median}{Median}: \omterm{prop:g10:coordgeom:midpoint}{midpoint}
of the $4$th and $5$th values, $\frac{8 + 9}{2} = 8.5$. \omterm{def:g11:stat:quartiles}{Quartiles}: a
quarter of $8$ is $2$, so $Q_1$ is the $2$nd value, $7$; three quarters is
$6$, so $Q_3$ is the $6$th value, $10$. \omterm{def:g11:stat:quartiles}{Interquartile range}:
$10 - 7 = 3$.
\end{solution}

\begin{solution}{exo:g11:stat:2}
Sum $40$, so $\bar x = 5$. Sum of squares:
$4 + 16 + 16 + 16 + 25 + 25 + 49 + 81 = 232$. By the shortcut formula,
\[
V = \frac{232}{8} - 5^2 = 29 - 25 = 4,
\qquad
\sigma = 2 .
\]
\end{solution}

\begin{solution}{exo:g11:stat:3}
\omterm{def:g11:stat:mean}{Mean}:
\[
\bar x = \frac{1 \times 6 + 2 \times 9 + 3 \times 8 + 4 \times 10
+ 5 \times 9 + 6 \times 8}{50}
= \frac{181}{50} = 3.62 .
\]
\omterm{def:g11:stat:median}{Median}: the cumulative frequencies are $6, 15, 23, 33, \dots$, so the
$25$th and $26$th sorted outcomes are both $4$: the \omterm{def:g11:stat:median}{median} is $4$.
\end{solution}

\begin{solution}{exo:g11:stat:4}
The total of the $15$ marks is $15 \times 11.2 = 168$; with the
newcomer, $\frac{168 + 18}{16} = \frac{186}{16} = 11.625$.
\end{solution}

\begin{solution}{exo:g11:stat:5}
$\frac{6 + 8 + x + 12 + 14}{5} = 10$ gives $40 + x = 50$, so $x = 10$.
The deviations of $6, 8, 10, 12, 14$ from the \omterm{def:g11:stat:mean}{mean} $10$ are
$-4, -2, 0, 2, 4$, with squares summing to $40$:
$V = \frac{40}{5} = 8$ and $\sigma = 2\sqrt2 \approx 2.83$.
\end{solution}

\begin{solution}{exo:g11:stat:6}
By \cref{prop:g11:stat:affine} with $a = \frac12$, $b = 4$: new \omterm{def:g11:stat:mean}{mean}
$\frac{12.4}{2} + 4 = 10.2$; new \omterm{def:g11:stat:variance}{standard deviation}
$\frac{2.5}{2} = 1.25$.
\end{solution}

\begin{solution}{exo:g11:stat:7}
The pooled \omterm{def:g11:stat:mean}{mean} is the \emph{weighted} average
\[
\frac{20 \times 10 + 30 \times 15}{50} = \frac{650}{50} = 13,
\]
not $12.5$: group B contributes more values, so it pulls the pooled \omterm{def:g11:stat:mean}{mean}
towards its own.
\end{solution}

\begin{solution}{exo:g11:stat:8}
Both lines are nearly on target on average ($1.002$ and $1.001$ kg), but
line 2's \omterm{def:g11:stat:variance}{standard deviation} is three times larger: its bags vary much
more, producing far more noticeably under- and over-filled bags. Line 2
should be dealt with first --- consistency, not just the \omterm{def:g11:stat:mean}{mean}, is what
needs fixing.
\end{solution}

\begin{solution}{exo:g11:stat:9}
$\bar x = \frac{70}{10} = 7$;
$V = \frac{540}{10} - 7^2 = 54 - 49 = 5$;
$\sigma = \sqrt5 \approx 2.24$.
\end{solution}

\begin{solution}{exo:g11:stat:10}
\emph{1.} With $23$: sum $50$, \omterm{def:g11:stat:mean}{mean} $5$; \omterm{def:g11:stat:median}{median} = \omterm{prop:g10:coordgeom:midpoint}{midpoint} of $5$th and
$6$th values $= 3$. Without $23$: sum $27$, \omterm{def:g11:stat:mean}{mean} $3$; \omterm{def:g11:stat:median}{median} (5th of 9
values) $= 3$.

\emph{2.} With $23$ ($n = 10$): $Q_1$ is the $3$rd value ($2$), $Q_3$ the
$8$th ($4$): \omterm{def:g11:stat:quartiles}{interquartile range} $2$. Without ($n = 9$): $Q_1$ is the
$3$rd value ($2$), $Q_3$ the $7$th ($4$): still $2$.

\emph{3.} The single outlier drags the \omterm{def:g11:stat:mean}{mean} from $3$ up to $5$ but leaves
the \omterm{def:g11:stat:median}{median} and the \omterm{def:g11:stat:quartiles}{interquartile range} unchanged: those two are
\emph{robust}. When a dataset may contain aberrant values, report the
\omterm{def:g11:stat:median}{median} and \omterm{def:g11:stat:quartiles}{interquartile range} rather than the \omterm{def:g11:stat:mean}{mean} and \omterm{def:g11:stat:variance}{standard
deviation}.
\end{solution}
